\section{Normalization and the underlying Einstein existence problem}
\label{sec:normalized-existence}

The integral obstruction, an approximate metric construction, and an
Einstein repair theorem answer different quantified questions. We record
the elementary reductions that place them in the same geometric problem.
In this section the dimension is any fixed $n\ge2$, and manifolds are
closed, connected and smooth; orientation is required only when a
four-dimensional characteristic number is used.

Let $\mathscr E(M)$ denote the set of all pairs $(h,\lambda)$ with
$h$ a smooth Riemannian metric on $M$ and $\Ric_h=\lambda h$.
For $\sigma\in\{-1,1\}$ put
\[
 \delta_\sigma(g)=
 \sup_{x\in M}\|g^{-1}\Ric_g-\sigma\Id\|_{\mathrm{op}}.
\]
Write $P_n^\sigma$ for the assertion that there exists
$\varepsilon_n^\sigma\in(0,1)$ such that, for every $M^n$ and every
$g$, the inequality $\delta_\sigma(g)\le\varepsilon_n^\sigma$ implies
$\mathscr E(M)\ne\varnothing$. The output metric need not lie near $g$
and its Einstein constant need not have sign $\sigma$.
Write $Q_n$ for the assertion that there exist real numbers $a<b$
such that, for every $M^n$ and $g$,
\[
 ag\le\Ric_g\le bg\quad\Longrightarrow\quad
 \mathscr E(M)\ne\varnothing.
\]
The inequalities are inequalities of quadratic forms, not componentwise
coordinate inequalities or scalar-curvature bounds.

\begin{proposition}[The two normalized branches]
\label{prop:normalized-branches}
In the same manifold category,
\[
 Q_n\quad\Longleftrightarrow\quad P_n^+\ \text{or}\ P_n^-.
\]
No sign assumption on the interval $[a,b]$ is needed for this equivalence.
\end{proposition}
\begin{proof}
For $c>0$, the covariant Ricci tensor satisfies
$\Ric_{cg}=\Ric_g$, whereas
$(cg)^{-1}\Ric_{cg}=c^{-1}g^{-1}\Ric_g$.
If $P_n^\sigma$ holds, take
$[a,b]=[\sigma-\varepsilon_n^\sigma,
\sigma+\varepsilon_n^\sigma]$.
Conversely, every nondegenerate interval $[a,b]$ contains a
nondegenerate closed interval $[u,v]$ with $uv>0$.
Set
\[
 \mu=\frac{u+v}{2},\qquad
 \sigma=\operatorname{sign}\mu,\qquad
 \varepsilon=\frac{v-u}{|u+v|}\in(0,1).
\]
If $\delta_\sigma(h)\le\varepsilon$, then $g=h/|\mu|$ has all
Ricci endomorphism eigenvalues in $[u,v]\subset[a,b]$.
The assertion $Q_n$ therefore supplies an Einstein metric on $M$.
\end{proof}

For a fixed sign define
\[
 \Gamma_n^\sigma=
 \inf_{\mathscr E(M^n)=\varnothing}\ \inf_{g\in\Met^\infty(M)}
       \delta_\sigma(g),
 \qquad \inf\varnothing=+\infty.
\]
Then $P_n^\sigma$ is equivalent to $\Gamma_n^\sigma>0$.
Indeed a positive infimum supplies a smaller positive threshold;
conversely a threshold bounds the infimum from below. Thus the negation
of one branch requires metrics $g_j$ on Einstein-free manifolds $M_j$
of one fixed dimension with $\delta_\sigma(g_j)\to0$.
The manifolds may vary with $j$. A fixed-manifold construction is
stronger than required, but it must still exclude every Einstein output.
The negation of $Q_n$ requires failure of both signs in that same
dimension; their underlying manifolds may be different.

In dimension four, taking $\lambda=\sigma$ in the definition of
$\E$ and comparing the tensor and operator norms gives
\begin{equation}\label{eq:normalized-pointwise-integral}
 \E(g)\le\fint_M|\Ric_g-\sigma g|^2
          \le4\delta_\sigma(g)^2.
\end{equation}
Consequently an all-metric estimate $\E\ge\eta>0$ on a specified
actual-input family excludes pointwise approximations below
$\sqrt\eta/2$ in that family. The converse implication is false:
Section~\ref{sec:integral-limits} constructs Einstein-free manifolds
with $\E\to0$ while the pointwise defect does not vanish.
Likewise, excluding an Einstein metric in a prescribed repair
neighborhood does not assert $\mathscr E(M)=\varnothing$.

These distinctions explain the roles of the constructions and global
topology in the supplements. They test the geometric limits of the
integral theory and the possibility of exact realization; they do not
silently replace its actual-input hypotheses by a dimension-only claim.
